\SourceLeaf{337}

Die Zeichen von
\[ \sin(\mathfrak{B}-B), \qquad \sin(\mathfrak{C}-C), \qquad \sin(\mathfrak{B}+\mathfrak{C}-B-C) \]
werden in den acht möglichen Combinationen diese sein

\begin{center}
\begin{tabular}{cc|cc}
I 1 & $+\;+\;+$ & II 1 & $-\;-\;-$ \\
I 2 & $-\;+\;-$ & II 2 & $+\;-\;+$ \\
I 3 & $+\;-\;-$ & II 3 & $-\;+\;+$ \\
I 4 & $-\;-\;+$ & II 4 & $+\;+\;-$ \\
\end{tabular}
\end{center}

Zur physischen Möglichkeit der Aufgabe ist also erforderlich und zureichend, dass die drei Sinus gleiche Zeichen haben.

Hat der erste Sinus ein von den beiden andern abweichendes Zeichen, so (Fall 3) gibt die Construction einen Punkt, der im Kreise $\gamma$ im ungeeigneten Segmente liegt.

Hat der zweite Sinus ein von den beiden andern verschiedenes Zeichen, so [Fall 2] liegt der durch die Construction gefundene Punkt im ungeeigneten Segmente des Kreises $\beta$.

\Needspace{6\baselineskip}Weicht hingegen der dritte Sinus rücksichtlich des Zeichens von den beiden andern ab, so [Fall 4] liegt der durch die Construction gefundene Punkt in beiden Kreisen $\beta$ und $\gamma$ in ungeeigneten Segmenten.

\hfill 1852 Juli 29.

\medskip

\begin{center}
\textbf{BEMERKUNGEN.}
\end{center}

Die meisten der vorstehenden Notizen sind in einem Handbuche enthalten, nur [7] befindet sich auf einem einzelnen Blatte. Was die Zeit der Abfassung betrifft, so stammt die Notiz [1] aus dem Jahre 1832, während [5] und [6] in das Jahr 1849, [7] in das Jahr 1852 zu setzen ist.

Vom December 1839 bis Ostern 1840 hat \textsc{Gauss} eine Vorlesung über die \emph{Theorie der imaginären Grössen} gehalten, die E. \textsc{Heine} gehört und ausgearbeitet hat. Dieser Ausarbeitung sind die folgenden auf das Viereck bezüglichen Ausführungen entnommen, die zur Erläuterung der unter [III] zusammengestellten Notizen dienen mögen. \hfill \textsc{Stäckel}.
